\documentclass[11pt,a4paper]{article}

\usepackage[margin=1in]{geometry}
\usepackage{amsmath,amssymb,amsthm}
\usepackage{enumitem}
\usepackage{hyperref}
\usepackage{xcolor}
\usepackage{titlesec}
\usepackage{fancyhdr}
\usepackage{booktabs}

\newtheorem{definition}{Definition}[section]
\newtheorem{remark}{Remark}[section]

\titleformat{\section}{\large\bfseries\color{blue!70!black}}{}{0em}{}[\titlerule]
\titleformat{\subsection}{\normalsize\bfseries\color{blue!50!black}}{}{0em}{}

\pagestyle{fancy}
\fancyhf{}
\fancyhead[L]{\small\textit{SAM3-DualZero-Conoid Foundations}}
\fancyhead[R]{\small Core Definitions (Part I)}
\fancyfoot[C]{\thepage}
\renewcommand{\headrulewidth}{0.4pt}

\title{\textbf{Foundations of the SAM3-DualZero-Conoid Model}\\[0.3em]
\large Part I: Right Conoid Geometry and Dual-Zero Hyperreal Structure}
\author{Shawn Dykes (\texttt{@Pieces1121}) \\ \small In collaboration with Grok}
\date{\today}

\begin{document}

\maketitle

\begin{abstract}
This document collects the core geometric and analytic formulas for the SAM3-DualZero-Conoid framework. It includes both rigorously derived expressions and key claimed relations from the model. Gaps are explicitly noted.
\end{abstract}

\section{Right Conoid Geometry}

\subsection{Definition and Metric}

The infinite right conoid is parametrized by
\[
\mathbf{r}(u,v) = \bigl( v \cos u,\ v \sin u,\ h(u) \bigr).
\]

For the **right helicoid** (\(h(u) = c u\)) the induced metric is
\begin{equation}
ds^2 = (v^2 + c^2)\, du^2 + dv^2.
\end{equation}

\subsection{Curvature}

Gaussian curvature:
\begin{equation}
K = -\frac{c^2}{(v^2 + c^2)^2} < 0.
\end{equation}

Mean curvature: \( H = 0 \) (minimal surface).

\subsection{Orthonormal Coframe and Volume Form}

\[
e^1 = \sqrt{v^2 + c^2}\, du, \quad e^2 = dv,
\qquad
\mathrm{vol} = \sqrt{v^2 + c^2}\, du \wedge dv.
\]

\subsection{Dirac Operator (Schematic)}

\[
D = i \gamma^1 \left( \frac{1}{\sqrt{v^2 + c^2}} \partial_u + \cdots \right) 
  + i \gamma^2 \partial_v + \text{connection terms}.
\]

\begin{remark}
The spectrum of \(D\) on the infinite helicoid is continuous. Dual-Zero regularization is required to obtain a well-defined spectral action.
\end{remark}

\section{Dual-Zero Hyperreal Structure}

\subsection{Definition}

\begin{definition}[Dual-Zero Hyperreal Structure]
A Dual-Zero structure on the hyperreals \(*\mathbb{R}\) consists of a pair of signed infinitesimals \(\varepsilon, \delta\) together with a duality involution \(*\) satisfying \(\varepsilon^* = \delta\) and \(\delta^* = \varepsilon\).
\end{definition}

\subsection{Proposed Scaling with Parameters}

The following relations are used in the model:
\begin{equation}
\varepsilon = \ell_0 \exp(-1/\omega_0), \qquad 
\delta = -\ell_0 \exp(-\omega_0).
\end{equation}

Here:
- \(\ell_0 > 0\) is a fundamental length scale,
- \(\omega_0 \approx 0.97\) is the Dual-Zero regulator strength.

\subsection{Claimed Physical Formulas}

\textbf{Newton’s Constant} (claimed from spectral action):
\begin{equation}
G_N = \frac{64\pi \ell_0^2}{45}.
\end{equation}

Other claimed relations (from spectral action expansion):
- Higgs boson mass \(\approx 126.2 \pm 2.05\) GeV (tuned via \(\omega_0\)),
- Sum of neutrino masses \(\sum m_\nu \approx 0.0585\) eV,
- Gauge coupling unification near \(10^{15.8}\) GeV.

\begin{remark}
The explicit derivation of \(G_N\) and the mass predictions from the spectral action has not yet been shown in detail in these foundational sections.
\end{remark}

\subsection{Interface with Binary Icosahedral Symmetry (Schematic)}

When Dual-Zero acts on representations \(\rho\) of the binary icosahedral group \(2I\), the model proposes a multiplicity-generating term of the form:
\[
\varepsilon \cdot \rho(g) + \delta \cdot \rho(g)^*.
\]

This is claimed to produce exactly **three chiral generations**.

\begin{remark}
The precise mechanism by which the above expression yields three generations is not yet derived in detail.
\end{remark}

\section{Critical Gaps (Summary)}

\begin{remark}
The following elements still require explicit formulas or derivations:
\begin{itemize}
    \item Precise algebraic rules and commutation relations for \(\varepsilon\) and \(\delta\).
    \item How Dual-Zero modifies the cutoff function or trace in the spectral action.
    \item Full derivation of \(G_N = \frac{64\pi \ell_0^2}{45}\).
    \item Explicit map from Dual-Zero + \(2I\) to three fermion generations.
    \item Concrete computational example of Dual-Zero regularization.
\end{itemize}
\end{remark}

\end{document}
